\section{Memory transactions and semantic validity}\label{sec:C17:transactions}
Transactional consistency preserves relationships among reads and writes; semantic validity concerns whether the recorded assertions and procedures are appropriate. A database can provide atomic updates that commit a false summary perfectly. A language model can produce an accurate summary whose references become inconsistent because another writer changes the source before commit. An agent memory system needs both a storage contract and a semantic admission contract.

Let a proposed update read objects $R$ at versions $v_R$ and produce writes $W$. An optimistic commit validates that the relevant read versions still match. If not, the update should be recomputed, merged under a justified rule, or rejected. The hard part is identifying the relevant read set. Generated abstractions can depend on a source that was seen earlier but omitted from the recorded extraction input. Missing provenance defeats conflict detection even if the database transaction implementation is correct.

The RecMem consolidation path in Chapter~\ref{ch:C14} illustrates a multi-layer publication boundary: episodes, semantic facts, and raw-memory deletion are distinct operations~\cite{repo_recmem2026}. HyperSkill maintenance similarly changes nodes, embeddings, and membership relations~\cite{repo_hyperskill2026}. A partially persisted update can make one representation disagree with another. The required invariant is not simply ``all writes completed,'' but that every published object has valid dependencies and can be interpreted under the same versioned schema.

Semantic validity needs richer conflict classes. A later observation can supersede a current fact without invalidating the truth of its historical episode. A skill can become incompatible with a new tool while remaining an accurate record of a past successful procedure. An authorization can expire without deleting the evidence that it once existed. These cases require status transitions, not indiscriminate overwrite or deletion.

Linearizable memory operations can support reasoning about a single shared object's update order~\cite{herlihy1990linearizability}. They do not determine whether the chosen semantic resolver is correct. A last-write-wins fact register can be linearizable while choosing the wrong assertion because the later writer had lower authority or described an older valid-time interval. Storage guarantees should therefore be stated using the actual abstract data type and application semantics.

A proposed semantic transaction writes candidate assertions with source links, validates role and authority, resolves temporal status, and atomically publishes a coherent view. If validation is model-based, retain its uncertainty and evidence rather than calling the transaction verified in the formal sense. For high-consequence obligations, deterministic fields can constrain what the model may revise. The goal is not to make every memory fact formally provable, but to prevent unexamined semantic changes from masquerading as ordinary storage maintenance.

